% ----------------------------------------------------------
\chapter{Estudo das alternativas de traçado}
% ----------------------------------------------------------

Os traçados foram escolhidos de forma a respeitar os requisitos mínimos e máximos de DNER (1999). 


\section{Traçado A}

O \textbf{Traçado A} possui extensão total de \textbf{5 Km e 701,48 m}, apresentando \textbf{4 curvas horizontais} e \textbf{8 curvas verticais}. O desenvolvimento do perfil horizontal está ilustrado na Figura~\ref{fig:P_H_A}, enquanto o perfil vertical é apresentado na Figura~\ref{fig:P_V_A}.

\begin{figure}[h]
    
    \includegraphics[width=0.85\textwidth]{unbtex-example/figuras/8.Alinhamento_A_horizontal.pdf}
    \caption{Perfil horizontal do Traçado A (1:200). (Todos os mapas estão orientados de acordo com o norte da quadrícula.)}
    \label{fig:P_H_A}
\end{figure}

\begin{figure}[h]
    
    \includegraphics[width=0.85\textwidth]{unbtex-example/figuras/9.alinhamento_1_perfil_total.pdf}
    \caption{Perfil vertical do Traçado A (h: 1:2000 e v: 1:200).}
    \label{fig:P_V_A}
\end{figure}

Em planta, o Traçado A é composto por \textbf{cinco tangentes}, com extensões de \textbf{967 m}, \textbf{2,5 Km}, \textbf{745 m}, \textbf{703 m} e \textbf{720 m}, respectivamente, partindo do Ponto de Partida (PP) até o Ponto Final (PF).

Em perfil, o traçado apresenta \textbf{oito curvas verticais} e \textbf{nove rampas}, com comprimentos de \textbf{652,03 m}, \textbf{302,90 m}, \textbf{161,73 m}, \textbf{397,73 m}, \textbf{496,36 m}, \textbf{870,90 m}, \textbf{1210,44 m}, \textbf{674,56 m} e \textbf{736,48 m}, respectivamente, de PP até PF. As inclinações correspondentes a essas rampas são, respectivamente, de \textbf{3,73\%}, \textbf{3,88\%}, \textbf{2,16\%}, \textbf{0,09\%}, \textbf{1,01\%}, \textbf{3,96\%}, \textbf{0,41\%}, \textbf{0,75\%} e \textbf{0,05\%}.
Além desses parâmetros, esse traçado possui deflexão máxima de 34°.













\section{Traçado B}

O \textbf{Traçado B} possui extensão total de \textbf{5 Km e 18,94 m}, apresentando \textbf{3 curvas horizontais} e \textbf{9 curvas verticais}. O desenvolvimento do perfil horizontal está ilustrado na Figura~\ref{fig:P_H_B}, enquanto o perfil vertical é apresentado na Figura~\ref{fig:P_V_B}.

\begin{figure}[h]
    
    \includegraphics[width=0.85\textwidth]{unbtex-example/figuras/10.Alinhamento_B_horizontal.pdf}
    \caption{Perfil horizontal do Traçado B (1:200). (Todos os mapas estão orientados de acordo com o norte da quadrícula.)}
    \label{fig:P_H_B}
\end{figure}

\begin{figure}[h]
    
    \includegraphics[width=0.85\textwidth]{unbtex-example/figuras/11.alinhamento_2_perfil_total.pdf}
    \caption{Perfil vertical do Traçado B (h: 1:2000 e v: 1:200).}
    \label{fig:P_V_B}
\end{figure}

Em planta, o Traçado B é composto por \textbf{cinco tangentes}, com extensões de \textbf{900,76 m}, \textbf{1275,03 m}, \textbf{1053,01 m} e \textbf{1790,13 m}, respectivamente, partindo do Ponto de Partida (PP) até o Ponto Final (PF).

Em perfil, o traçado apresenta \textbf{sete curvas verticais} e \textbf{oito rampas}, com comprimentos de \textbf{422,99 m}, \textbf{331,56 m}, \textbf{685,29 m}, \textbf{37,81 m}, \textbf{567,81 m}, \textbf{551,55 m}, \textbf{631,79 m}, \textbf{584,84 m}, \textbf{376,18 m} e \textbf{829,10 m}, respectivamente, de PP até PF. As inclinações correspondentes a essas rampas são, respectivamente, de \textbf{4,48\%}, \textbf{3,39\%}, \textbf{4,48\%}, \textbf{1,00\%}, \textbf{3,85\%}, \textbf{3,04\%}, \textbf{2,82\%}, \textbf{2,45\%}, \textbf{3,46\%} e \textbf{4,39\%}.
Além desses parâmetros, esse traçado possui deflexão máxima de 32°.